% ---------------------------------------------------------------------------
% Section 11: supplements to Sections 6, 7 and 9 (p = 2, n = 2 and n = 3).
% Placement in the source of the paper:
%   11.1 -> after Theorem 6.2 (it settles the "[open]" item there and the p = 2
%           part of Question 6.4);
%   11.2 -> after Proposition 3.7 / Remark 4.5 (Riemann-Roch bound), and its
%           corollary at (2,3) after Theorem 7.1;
%   11.3 -> after Remark 7.6 (Case A empty; U_3 never a quadratic witness; new
%           families with ed = 2; structure of a witness);
%   11.4 -> after Proposition 7.5 / in Section 9 (the cubic slot);
%   11.5 -> replaces the status table of Theorem 7.1 and updates Section 9.2.
% Appendices A-E hold the supporting material.
% ---------------------------------------------------------------------------
\section{Supplements at \texorpdfstring{$(p,n)=(2,2)$ and $(2,3)$}{(p,n)=(2,2) and (2,3)}:
the cubic slot, quadratic witnesses, and a Riemann--Roch bound}\label{sec:supp}

\noindent\textit{Addendum (27 September 2026).} This section settles the open item
of Theorem~6.2, gives a general Riemann--Roch bound for the prime-to-$p$ jump degree,
shows that the Artin--Schreier--Witt curve never carries a quadratic witness at $(2,3)$
and that a quadratic witness can never have a descent field inside $K$, produces new
families of quadratic extensions with essential dimension exactly $2$, and reformulates
the cubic slot. Notation is that of Sections 2, 5--7: $k$ algebraically closed of
characteristic $2$, $K=k(s_0,s_1,s_2)$, $L=k(z_0,z_1,z_2)$, $s=\wp(z)$, $K_2=k(s_0,s_1)$,
$\tau_2$ the generic $\Z/4$-torsor, $E'\colon w^2+w=u(u+1)(u+s_0)+s_1$, $\bar U_n$ the
Artin--Schreier--Witt curve $\wp y=[t]$, $\Phi_2$ the polynomial of Proposition~7.5 (written
out in \eqref{eq:Phi2}). We write $(s_0,s_1)_K$ for the $\Z/4$-class $(s_0,s_1)\in W_2(K)/\wp$,
i.e.\ $\tau_2\otimes_{K_2}K$; by Lemma~4.1(i), $\ed_k((s_0,s_1)_K)=\ed_k(\tau_2)=2$. Every
separable quadratic extension of $K$ is $K(v_f)$, $v_f^2+v_f=f$, $f\in K/\wp K$
(Proposition~3.6). Scripts referred to are listed in Appendix~\ref{app:scripts}.

\subsection{Supplement to Theorem 6.2: the prime-to-\texorpdfstring{$2$}{2} jump degree at
\texorpdfstring{$n=2$}{n=2}}

\begin{theorem}\label{thm:A}
Let $p=2$, $n=2$. Let $u$ be a root of
\[
P(u):=u^3+(1+s_0)u^2+s_0u+s_1\;=\;u(u+1)(u+s_0)+s_1 .
\]
Then $F'=K_2(u)$ is a separable cubic extension of $K_2$ and
$(s_0,s_1)\equiv[s_0+u^2+u]$ mod $\wp W_2(F')$. Hence $\ed_k(\tau_2\otimes F')=1$ and
\[
d_1^{(p')}(\tau_2)=3 .
\]
More precisely, the cubic closed points of $E'$ over $K_2$ are exactly the sections by
the lines $w=\alpha u+\beta$, $\alpha,\beta\in K_2$; each is a single free point of degree
$3$, and over its residue field $(s_0,s_1)$ is Teichm\"uller. Consequently
$D(U_2)=\Z_{\ge2}$, and the profile of $\tau_2$ is completely known:
$\edd1=2$, $\edd2=\edd3=1$, $\edd4=0$, $d_1=2$, $d_1^{(p')}=3$, $d_0=4$.
\stag{proved}
\end{theorem}

\begin{proof}
$P\in k(s_0)[u][s_1]$ is of degree $1$ in $s_1$ with coprime coefficients $1$ and
$u(u+1)(u+s_0)$, hence irreducible in $k(s_0)[u,s_1]$, hence (Gauss, $P$ monic in $u$)
irreducible in $K_2[u]$; an irreducible cubic in characteristic $2$ is separable. By
(3)--(4), for $h=(u,0)\in W_2(F')$: $(s+\wp h)_0=s_0+u^2+u$ and
$(s+\wp h)_1=s_1+u^3+u^2+s_0u^2+s_0u=P(u)=0$, so $s+\wp h=[s_0+u^2+u]$; equivalently
$(u,0)$ is an $F'$-point of $E'^\circ=\tau_2U_2$ and the last sentence of Theorem~6.2
applies. So $\ed_k(\tau_2\otimes F')\le1$, and $\ge1$ since $3<4=d_0$. As
$d_1^{(p')}$ is odd and $\ge d_1(\tau_2)=2$ (Proposition~6.1), $d_1^{(p')}(\tau_2)=3$.

Cubic points: $E'(K_2)=\{\infty'\}$ (Theorem~6.2), so $\operatorname{Pic}^0(E')(K_2)=0$
and every effective $K_2$-divisor of degree $3$ is linearly equivalent to $3\infty'$,
i.e.\ is the zero divisor of an element of $L(3\infty')=\langle1,u,w\rangle$; the
sections $u=c$ are a degree-$2$ fibre plus $\infty'$, so the cubic points are the
sections by lines $w=\alpha u+\beta$. Such a section is a single closed point of degree
$3$ because a $K_2$-rational component would be a $K_2$-point of $E'^\circ$. It lies in
$E'^\circ$, hence is free. $D(U_2)\supseteq\{2,3\}$ and Appendix~\ref{app:genus} gives
all of $\Z_{\ge2}$.
\end{proof}

\begin{remark}
(i) This closes the item marked [open] in Theorem~6.2 and the $p=2$ part of
Question~6.4: the ``group-theoretic'' obstruction $E'(K_2)=0$ kills degree-$1$ points and
\emph{produces} the cubic ones. (ii) For odd $p$ the same mechanism gives
$d_1^{(p')}(\tau_2)\le p^2-p+1$ (Theorem~\ref{thm:B}); Question~6.4 for degrees in
$[2,p-1]$ is untouched.
\end{remark}

\subsection{Supplement to Proposition 3.7 and Remark 4.5: a Riemann--Roch bound}

Let $l_n=l_n(p)$ be the largest lower ramification break of $\bar U_n\to\mathbb P^1_t$ at
$\infty$, i.e.\ the pole order at $\infty$ of the reduced top Witt coordinate of $U_n$;
by Appendix~\ref{app:genus},
\begin{equation}\label{eq:ln}
l_n(p)=1+(p-1)\sum_{j=1}^{n-1}p^{2j-1}:\qquad
l_2(2)=3,\quad l_2(p)=p^2-p+1,\quad l_3(2)=11,\quad l_4(2)=43 .
\end{equation}

\begin{theorem}\label{thm:B}
For every $(p,n)$, $D(U_n)$ contains an integer prime to $p$ and $\le l_n$; hence
\[
d_1^{(p')}(\tau_\gen)\le l_n(p).
\]
This is the best bound obtainable from the pencils $|d\cdot\infty''|$ on
$\tau_\gen\bar U_n$: every function regular on the affine curve $U_n$ with pole order prime
to $p$ has pole order $\ge l_n$.
\stag{proved; the closed formula \eqref{eq:ln} is proved for $(2,2)$, $(2,3)$, $(p,2)$
and cited in general}
\end{theorem}

\begin{proof}
$X:=\tau_\gen\bar U_n$ is a smooth projective geometrically integral $K$-curve with the
$K$-point $\infty''$ (twist of the $G$-fixed point), and its free locus is
$X\smallsetminus\{\infty''\}$ (Proposition~3.7(c)). Riemann--Roch dimensions
$\ell(d\infty'')$ are invariant under base change to $\bar K$, where
$(X,\infty'')\cong(\bar U_n,\infty)$; so the Weierstrass semigroup of $\infty''$ on $X$ is
that of $\infty$ on $\bar U_n$. The top layer $U_n\to U_{n-1}$ is $y^p-y=c$ with
$v_\infty(c)=-l_n$ on $U_{n-1}$ and $\infty$ totally ramified, so $v(y)=-l_n$ on $U_n$
and $l_n$ is a pole order. Hence there is $f\in K(X)$ with pole divisor exactly
$l_n\infty''$. For $c\in K$ the divisor $\operatorname{div}(f-c)+l_n\infty''$ is
effective, $K$-rational, of degree $l_n$ and supported on $X\smallsetminus\{\infty''\}$.
Its closed points have degrees summing to $l_n\equiv1\pmod p$, so one of them has degree
$e$ prime to $p$, $e\le l_n$, and it is free; Theorem~4.2 gives $\edd e_k(\tau_\gen)\le1$.
For optimality: the coordinate ring of $U_n$ is free over that of $U_{n-1}$ with basis
$1,y,\dots,y^{p-1}$, and $v(a_jy^j)=p\,v_{U_{n-1}}(a_j)-jl_n$ are pairwise distinct
modulo $p$, so a function with pole order prime to $p$ has a dominant term with $j\ge1$
and pole order $\ge l_n$.
\end{proof}

\begin{corollary}\label{cor:B23}
At $(p,n)=(2,3)$: the fibres of the pole-order-$11$ pencil are free divisors of degree
$11$ without degree-$1$ components (Proposition~7.5), so each contains a free point of
odd degree in $\{3,5,7,9,11\}$, and $11\in D(U_3)$ exactly (Appendix~\ref{app:genus}).
Hence
\[
d_1^{(p')}(\tau_\gen)\in\{3,5,7,9,11\},\qquad \delta_{\rm free}(U_3)\in\{3,5,7,9,11\},
\]
the second by Corollary~\ref{cor:noquad} below. This refines ``$d_1^{(p')}\ge3$ finite''
in Theorem~7.1. \stag{proved}
\end{corollary}

\begin{remark}
The naive argument ``$\ell(d\infty'')\ge2$ for $d\ge g+1$, take $d$ odd'' does not
work: the fibres of $f\in L(d\infty'')$ have degree equal to the actual pole order of
$f$, which lies in the Weierstrass semigroup and may be even. At $(2,3)$, $g(\bar U_3)=7$
and the semigroup of $\infty$ is $\langle4,6,11\rangle$, whose gaps are
$\{1,2,3,5,7,9,13\}$: every function of pole order $\le10$ has even pole order.
\end{remark}

\subsection{Supplement to Section 7: quadratic witnesses}

Fix $f\in K\smallsetminus\wp K$ and $K'=K(v_f)$. A \emph{witness} is such a $K'$ with
$\ed_k(\tau_\gen\otimes K')=1$; a \emph{descent field} is a subfield $k\subset F_0\subset
K'$ of transcendence degree $1$ with a class $\beta\in W_3(F_0)/\wp$, $s\equiv\beta$ in
$W_3(K')/\wp$.

\begin{theorem}[no descent field inside $K$]\label{thm:caseA}
Let $K'=K(v_f)$ be a witness with descent field $F_0$ and class $\beta$.
\begin{enumerate}
\item[(a)] If $F_0\subset K$ then $F_0=k(t)$ is rational, $s\equiv\beta+\varepsilon(0,0,f)$
mod $\wp W_3(K)$ with $\varepsilon\in\{0,1\}$, and $\ed_k(\tau_\gen)\le\tr_kk(t,f)\le2$.
\item[(b)] The case (a) does not occur: $F_0\not\subset K$.
\item[(c)] The truncation of $\beta$ is a descent of $(s_0,s_1)_K\otimes K'$, so
$f\in\mathfrak F:=\{f\in K/\wp K:\ed_k((s_0,s_1)_K\otimes K(v_f))=1\}$.
\end{enumerate}
\stag{proved}
\end{theorem}

\begin{proof}
(a) $\ker(H^1(K,\Z/8)\to H^1(K',\Z/8))=\operatorname{Hom}(\Z/2,\Z/8)\cong\Z/2$ by
inflation--restriction, generated by the induced class $V^2[f]=(0,0,f)$ (Section~2.2).
Since $s$ and $\beta$ both lie in $W_3(K)/\wp$ and agree over $K'$, $s\equiv\beta+
\varepsilon(0,0,f)=(\beta_0,\beta_1,\beta_2+\varepsilon f)$ by (2); and a subfield of
transcendence degree $1$ of $k(s_0,s_1,s_2)$ is $k(t)$ (a curve dominated by $\A^3$ is
dominated by a line, Lemma~2.5(i)). (b) Truncate the congruence of (a) to $W_2$: the
term $(0,0,f)$ dies and $(s_0,s_1)\equiv(\beta_0,\beta_1)$ mod $\wp W_2(K)$ with
$\beta_0,\beta_1\in F_0$, so $\ed_k((s_0,s_1)_K)\le1$, contradicting
$\ed_k((s_0,s_1)_K)=2$. (c) Truncation commutes with $\wp$ and with base change.
\end{proof}

So a quadratic witness, if one exists, has a descent field \emph{not} contained in $K$
and can never be turned into a counterexample to Ledet's conjecture by this route; the
truncation argument is the one used in the proof of Proposition~7.5 (the step
``$\varepsilon=1\Rightarrow$ a $K$-point of $E'$''), applied to an arbitrary descent
class instead of $[t]$.

\begin{corollary}[$\bar U_3$ admits no quadratic witness]\label{cor:noquad}
For every quadratic $K'/K$, $\tau_\gen\bar U_3$ has no free $K'$-point: $s$ is not
Teichm\"uller over any quadratic extension of $K$, and $D(U_3)\cap\{1,2\}=\emptyset$.
In particular Algorithm~9.1 has nothing to find, and the ``points not coming from the
$s_2$-line'' of Remark~7.6 do not exist. \stag{proved}
\end{corollary}

\begin{proof}
By Proposition~7.5 a free $K'$-point needs $Q=(u,w_1)\in E'(K')$ with $u\notin K$. Let
$\sigma$ generate $\operatorname{Gal}(K'/K)$. $Q+\sigma Q$ is $\sigma$-invariant, hence
lies in $E'(K)=\{\infty'\}$ (Lemma~4.1(ii)); so $\sigma Q=-Q=(u,w_1+1)$, whence
$\sigma(u)=u$ and $u\in K$.
\end{proof}

The same argument at level $2$ shows that $E'(K(v_f))\ne\{\infty'\}$ iff $f\equiv
s_1+u(u+1)(u+s_0)$ mod $\wp K$ for some $u\in K$: the $E_0$-witnesses of Theorem~6.2
are exactly the $u$-fibres (Proposition~\ref{prop:FE0}).

\begin{proposition}[fixed-field trick]\label{prop:fft}
Let $K'=K(v_f)$ with $f\notin\{0,s_0\}+\wp K$, so that $L'=L\otimes_KK'=L(v_f)$ is a
field. Suppose there are a subfield $M\subset L'$ and $x\in L'$ with $L'=M(x)$ purely
transcendental, such that $M$ is fixed pointwise by some $1\ne h\in G$. Then
$\ed_k(\tau_\gen\otimes K')\ge2$. The same holds for $(s_0,s_1)_K\otimes K'$, with $L'$
replaced by $k(z_0,z_1,s_2)(v_f)$ and $G$ by $\Z/4$. \stag{proved}
\end{proposition}

\begin{proof}
A witness gives, by Corollary~4.4, a smooth projective curve $C$ with faithful
$G$-action and a $G$-equivariant dominant rational map from a model of $L'$; take the
model $B\times\A^1$ with $k(B)=M$. Lemma~7.2 for $B\times\A^1\to B$ gives: either a fibre
$\A^1$ maps nonconstantly to $C$, so $C\cong\mathbb P^1$ (Lemma~2.5(i)), impossible
(Lemma~2.5(ii)); or $k(C)\subset M\subset(L')^h$, so $h$ acts trivially on $C$,
contradicting faithfulness.
\end{proof}

\begin{theorem}[quadratic extensions that are not witnesses]\label{thm:fam}
$\ed_k(\tau_\gen\otimes K(v_f))\ge2$, i.e.\ $K(v_f)$ is not a witness, for every
$f\notin\wp K$ in each of the classes
\begin{enumerate}
\item[(a)] $f\in K_2+\wp K$ (this contains $K(v_1)$ and $K(z_0)$ of Theorem~7.3);
\item[(b)] $f\in s_2+K_2+\wp K$ (this contains $K(v_2)$, not covered by Theorem~7.3);
\item[(c)] $f\in k(g)+\wp K$ for any $g\in K$ with $K=K_2(g)$, e.g.\ $f\in k(s_2)$,
$f=1/(s_2+s_0s_1)$;
\item[(d)] $f\in k(s_0,s_1)[s_2]$ a polynomial in $s_2$ of degree $\le7$.
\end{enumerate}
In cases (a), (c) already $\ed_k((s_0,s_1)_K\otimes K(v_f))=2$, i.e.\ $f\notin\mathfrak F$.
Equality $\ed_k(\tau_\gen\otimes K(v_f))=2$ holds in case (b), and in case (a) for
$f\equiv s_0$ or $f\equiv s_1$ (Theorem~7.3); for the other $f$ the value is $2$ or
$\ed_k(\Z/8)$, and only the lower bound is claimed.
\stag{proved}
\end{theorem}

\begin{proof}
In case (b) the upper bound is Theorem~7.1 (the class $(s_0,s_1,e)$ has a descent of
transcendence degree $2$); in general $\ed_k(\tau_\gen\otimes K')\le\ed_k(\tau_\gen)$. (a) Replace $f$ by $e\in K_2$ with $f\equiv e$. If
$e\equiv s_0$, Proposition~3.9. Otherwise $M:=k(z_0,z_1,v_e)$ is algebraic over
$k(z_0,z_1)$, $L'=M(z_2)$, and $g^4\colon z\mapsto z+V^2(1)=(z_0,z_1,z_2+1)$ fixes $M$;
apply Proposition~\ref{prop:fft}. (b) Write $f=s_2+e+\wp h$, $e\in K_2$. Over
$K'=K(v_f)$, $s_2=\wp(v_f+h)+e$, so $s=(s_0,s_1,e)+\wp V^2(v_f+h)\equiv(s_0,s_1,e)$ by
(2). Now $K'=K_2(v_f)$ is purely transcendental over $K_2$, so by Lemma~4.1(i)
$\ed_k(\tau_\gen\otimes K')=\ed_k((s_0,s_1,e)_{K_2})\ge\ed_k((s_0,s_1)_{K_2})=2$.
(c) Replace $f$ by an element of $k(g)$. Put $F:=k(g,v_f)$; then $K'=F(s_0,s_1)$ with
$s_0,s_1$ algebraically independent over $F$, and the total space of $(s_0,s_1)_K\otimes
K'$ is $k(z_0,z_1,g,v_f)=F(z_0,z_1)$, a field on which $\Z/4$ acts fixing $F$; the
argument of Proposition~\ref{prop:fft} with the fibration $B\times\A^2\to B$ excludes a
faithful $\Z/4$-curve, so $\ed_k((s_0,s_1)_K\otimes K')=2$, and a descent of $s$ over
$K'$ would truncate to a descent of $(s_0,s_1)_K\otimes K'$. (d) If $f\in K_2+\wp K$ use
(a); otherwise Theorem~\ref{thm:struct}(b) gives a witness curve of $p$-rank $0$ and
genus $\le6$, contradicting Corollary~\ref{cor:g7}.
\end{proof}

\begin{theorem}[structure of a quadratic witness]\label{thm:struct}
Let $K'=K(v_f)$ be a witness with $f\notin K_2+\wp K$, $C$ its curve, $C'=C/G$. For
$a\in k^2$ let $D_{f,a}\colon y^2+y=f(a_0,a_1,s_2)$ (a curve over $k$ in the $s_2$-line)
and, for $b\in\wp_2^{-1}(a)\subset\A^2_z$, let $D_b$ be its pullback along
$z_2\mapsto s_2=z_2^2+z_2+C(b_0,b_1)$.
\begin{enumerate}
\item[(a)] For $a$ in a dense open subset of $\A^2(k)$: $D_{f,a}$ is integral and
dominates $C'$, and every irreducible component of $D_b$ dominates $C$. Hence
$g(C)\le g(\hat D_b)$ and $\sigma(C)\le\sigma(\hat D_b)$ ($\sigma$ the $p$-rank, $\hat D_b$
the normalisation of a component).
\item[(b)] If $f\in k(s_0,s_1)[s_2]$ has $s_2$-degree $m\ge1$, then $\sigma(C)=0$,
$C\to C'$ is totally ramified at exactly one point and \'etale elsewhere, $C'$ has
$p$-rank $0$ and genus $\le\lfloor(m-1)/2\rfloor$, and $7\le g(C)\le m-1$. In particular
$m\ge8$, and for $8\le m\le15$ one has $C'\cong\mathbb P^1$: $C$ is the
Artin--Schreier--Witt curve of a polynomial Witt vector $\beta\in W_3(k[t])$, and
$s\equiv\beta(t)$ over $K'$ with $t\in K'\smallsetminus K$.
\end{enumerate}
\stag{proved}
\end{theorem}

\begin{proof}
(a) Let $B_f\colon y^2+y=f(s_0,s_1,s_2)$ over $\A^3_s$, $k(B_f)=K'$. Its generic fibre
over $\A^2_{(s_0,s_1)}$ is the Artin--Schreier curve $y^2+y=f$ over $K_2$ in the
$s_2$-line, geometrically integral because $f\notin K_2+\wp K$ (a constant quadratic
subextension $K_2(v_e)\subset K'$ would force $f\equiv e$). Apply Lemma~7.2 to $B_f\to
\A^2_s$ and the dominant rational map $B_f\dashrightarrow C'$ given by $k(C')\subset K'$:
the alternative $k(C')\subset k(s_0,s_1)\subset K$ is excluded by
Theorem~\ref{thm:caseA}(b), so $D_{f,a}\dashrightarrow C'$ is dominant for general $a$.
Let $Y_f$ be the normalisation of $\A^3_z\times_{\A^3_s}B_f$, with function field $L'$;
the equivariant map $Y_f\dashrightarrow C$ composed with $C\to C'$ factors through
$Y_f\to B_f\dashrightarrow C'$, and the fibre of $Y_f$ over $a$ is
$\bigsqcup_{b\in\wp_2^{-1}(a)}D_b$, permuted transitively by $G$. Each $D_b\to D_{f,a}
\to C'$ is nonconstant, hence so is $D_b\to C$. Genus and $p$-rank do not increase under
dominant maps of curves. (b) $D_{f,a}$ is $y^2+y=$ (polynomial of degree $m$ in $s_2$),
with a single branch point, so $\sigma(D_{f,a})=0$ (Deuring--Shafarevich) and
$g\le\lfloor(m-1)/2\rfloor$; $D_b$ is $y^2+y=$ (polynomial of degree $2m$ in $z_2$),
again with one branch point, so $\sigma=0$ and, after Artin--Schreier reduction to odd
degree $\le2m-1$, $g(D_b)\le m-1$. By (a), $\sigma(C)=\sigma(C')=0$ and $g(C)\le m-1$;
Lemma~\ref{lem:prank0} and Corollary~\ref{cor:g7} give the ramification structure and
$g(C)\ge7$. If $g(C')=1$ then Riemann--Hurwitz gives $2g(C)-2\ge0+28$, so $m\ge16$; for
$m\le15$, $C'=\mathbb P^1$ and the cover is \'etale outside one point, so its class
lies in $H^1(\A^1,\Z/8)=W_3(k[t])/\wp$, with $\beta_0$ nonconstant since
$s_0\notin\wp K'$ (else $K'=K(z_0)$).
\end{proof}

\begin{corollary}[minimal genus]\label{cor:g7}
A smooth projective curve with faithful $\Z/8$-action and $p$-rank $0$ in
characteristic $2$ has genus $\ge7$, with equality iff $C/G\cong\mathbb P^1$ and the
upper ramification breaks at the unique ramified point are $(1,2,4)$; $\bar U_3$ is an
example, $g(\bar U_3)=7$. A faithful $\Z/4$-curve of $p$-rank $0$ has genus $\ge1$, with
equality iff it is $(E_0,\text{Witt action})$ of Theorem~6.2. \stag{proved}
\end{corollary}

\begin{proof}
Lemma~\ref{lem:prank0}: one totally ramified point, $C/G$ of $p$-rank $0$. The different
exponent at that point is $(u_1+1)+2(u_2+1)+4(u_3+1)$ in terms of the upper breaks, with
$u_1\ge1$ and $u_{i+1}\ge2u_i$ \cite{Ga,Th}, so $2g-2\ge-16+28=12$. For $\Z/4$:
$2g-2\ge-8+2+6=0$, and the genus-$1$ case is $(E_0,\text{Witt})$ by
Lemma~\ref{lem:Z4ss}.
\end{proof}

\begin{remark}
(i) Theorem~7.3's proofs for $K(v_1)$ and $K(v_1,v_2)$ are instances of the mechanism
of Theorem~\ref{thm:struct}; for $f\in K_2$ the fixed-field trick is shorter, but for
the quartic tower it does not apply directly ($g^4$ moves $y'=v_2+z_2$) and the argument
of Theorem~7.3 is needed. (ii) The only genus-$1$ faithful $\Z/8$-curves are ordinary
elliptic curves with $G$ acting by translation by a point of order $8$
(Appendix~\ref{app:RH}); by Theorem~\ref{thm:struct}(a) they can occur as witness curves
only if $f$ has, for general $a$, at least two poles on the $s_2$-line, and then
$\tau_\gen\otimes K'$ is the \'etale $\Z/8$-isogeny torsor of an ordinary elliptic curve
at a non-constant point (Appendix~\ref{app:isog}). This and polynomial $f$ of
$s_2$-degree $\ge8$ are the only remaining shapes of a quadratic witness.
\end{remark}

\subsection{Supplement to Proposition 7.5 and Question (Q2): the cubic slot}

Let $F'/K$ be cubic (separable without loss, Proposition~3.6). Then $L\otimes_KF'$ is a
field and Corollary~4.4 applies: a witness for $\edd3(\tau_\gen)=1$ is a faithful
$\Z/8$-curve $C$ with a free $F'$-point on $\tau_\gen C$.

\begin{proposition}\label{prop:cubiccurves}
A cubic witness curve $C$ has genus $\ge3$ and, if its $p$-rank is $0$, genus $\ge7$.
Ordinary elliptic curves with translation action carry no free point of odd degree.
\stag{proved}
\end{proposition}

\begin{proof}
$\mathbb P^1$ and supersingular elliptic curves are excluded by Lemma~2.5. If $C$ is an
ordinary elliptic curve with $G$ acting by translation by $P$ of order $8$ (the only
genus-$1$ case, Appendix~\ref{app:RH}), then $\tau_\gen C$ is a principal homogeneous
space under $C$ whose period divides $8$; period and index have the same prime divisors
\cite{LT}, so a nontrivial such torsor has no closed point of odd degree, and a trivial
one has a $K$-point, free since the action is free, forcing $\ed_k(\tau_\gen)\le1$,
false. Genus $2$: the hyperelliptic involution is central, and the image of $\Z/8$ in
$\operatorname{Aut}(\mathbb P^1)$ is $\Z/8$ or $\Z/4$, excluded by Lemma~2.5(ii).
$p$-rank $0$: Corollary~\ref{cor:g7}.
\end{proof}

On $\bar U_3$ the question is completely explicit. By Proposition~5.2(a) and (3)--(4), a
free $F'$-point of $\tau_\gen\bar U_3$ is $h=(u,w_1,w_2)\in W_3(F')$ with
\begin{equation}\label{eq:U3eq}
w_1^2+w_1=u(u+1)(u+s_0)+s_1,\qquad w_2^2+w_2=e_Q:=s_2+\Phi_2(s_0,s_1,u,w_1),
\end{equation}
and then $s\equiv[s_0+u^2+u]$ over $F'$.

\begin{proposition}[reformulation]\label{prop:C}
$\edd3(\tau_\gen)=1$ is witnessed on $\bar U_3$ iff there are $\alpha,\beta\in K$ such
that, for $u$ a root of
\[
u^3+(1+s_0+\alpha^2)u^2+(s_0+\alpha)u+(s_1+\beta^2+\beta)=0,
\]
one has $e_Q=s_2+\Phi_2(s_0,s_1,u,\alpha u+\beta)\in\wp K(u)$. Every such cubic is
irreducible over $K$, and these are all the cubic points of $E'$ over $K$.
\stag{proved}
\end{proposition}

\begin{proof}
As in Theorem~\ref{thm:A}, $E'(K)=\{\infty'\}$ (Lemma~4.1(ii)) gives
$\operatorname{Pic}^0(E')(K)=0$, so the cubic points of $E'$ over $K$ are the line
sections, each irreducible. A free $F'$-point with $[F':K]=3$ has $F'=K(u)$: if $u\in K$
then $t=s_0+u^2+u\in K$ and $s\equiv[t]$ over $F'$ would give $s\equiv[t]$ over $K$,
since $\res_{F'/K}$ is injective on $H^1(\cdot,\Z/8)$ ($\cores\circ\res=3$ is a unit
mod $8$).
\end{proof}

\begin{proposition}[necessary condition and a partial obstruction]\label{prop:D}
\begin{enumerate}
\item[(a)] If $s\equiv[t]$ over a cubic $F'=K(u)$ then $s\equiv3\cdot\cores_{F'/K}[t]$ in
$W_3(K)/\wp$, where $\cores$ is the Witt-vector trace
$[t]+[t']+[t'']=(e_1,e_2,e_1e_3+e_1^2e_2)$ ($e_i$ the elementary symmetric functions of
the conjugates). Coordinates $0$ and $1$ of this condition hold automatically for every
line, and coordinate $2$ reads
\begin{equation}\label{eq:N1}
\operatorname{Tr}_{K(u)/K}(e_Q)=s_2+R(s_0,s_1,\alpha,\beta)\in\wp K
\end{equation}
with $R$ the explicit polynomial of Appendix~\ref{app:cubic}.
\item[(b)] \eqref{eq:N1} fails for all $\alpha,\beta\in K_2$, and more generally for all
$\alpha,\beta$ with pole order $\le1$ at $s_2=\infty$. Hence a cubic witness on
$\bar U_3$, if any, has a line with $\alpha$ or $\beta$ of pole order $\ge2$ at
$s_2=\infty$.
\end{enumerate}
\stag{(a) proved, coordinates checked symbolically; (b) proved, with computer-verified
coefficients}
\end{proposition}

\begin{proof}
(a) $\cores\circ\res=3$ on $H^1(K,\Z/8)$ and $3^2\equiv1\pmod8$. Corestriction is a
morphism of $\delta$-functors for $0\to\Z/8\to W_3\xrightarrow{\wp}W_3\to0$, and on
$H^0$ it is the norm, i.e.\ the Witt trace; the formula for $[t]+[t']+[t'']$ follows from
(3), and the coordinate computations are in Appendix~\ref{app:cubic}. (b) For
$\alpha,\beta\in K_2$, $s_2+R$ has a simple pole at $s_2=\infty$ (Remark~7.4). In
general write $\alpha=\sum_{i\le1}a_is_2^i$, $\beta=\sum_{i\le1}b_is_2^i$ in
$K_2((1/s_2))$. Membership in $\wp$ of the completion forces, for every odd $i_0\ge1$, the
chain condition $\sum_{j}c_{i_02^j}^{2^{J-j}}=0$ on the coefficients $c_i$ of $s_2^i$ in
$s_2+R$ (Appendix~\ref{app:cubic}). The only monomial of $R$ producing $s_2^{11}$ from
exponents $\le1$ is $\alpha^{11}$, giving $c_{11}=a_1^{11}$ and $a_1=0$; then
$c_3=b_1^3$ gives $b_1=0$; then $\alpha,\beta$ are integral at $\infty$ and
$s_2+R$ has a simple pole, contradicting the chain condition for $i_0=1$.
\end{proof}

Exhaustive searches over $\mathbb F_2$-polynomial lines (about $2.4\cdot10^6$ pairs
$(\alpha,\beta)$ of small degree, Appendix~\ref{app:cubic}) found no line satisfying
\eqref{eq:N1}. Note that \eqref{eq:N1} is only necessary: the kernel of $\cores$ on
$H^1(K(u),\Z/2)$ is large. Our expectation is that no cubic witness exists on
$\bar U_3$; whether another faithful $\Z/8$-curve of genus $\ge3$ carries one is a
separate question, so (Q2) remains open.

\subsection{Updated status at \texorpdfstring{$(2,3)$}{(2,3)} and revised questions}

\begin{theorem}[status at $(2,3)$]\label{thm:status}
\[
\edd1\in\{2,3\},\quad\edd2\in\{1,2\},\quad\edd3\in\{1,2\},\quad
\edd4=\dots=\edd7=1,\quad\edd8=0;
\]
$d_2\in\{1,2\}$, $d_1\in\{2,3,4\}$, $d_0=8$, $d_1^{(p')}\in\{3,5,7,9,11\}$. Moreover:
a quadratic witness for $d_1=2$ has no descent field inside $K$, is not $\bar U_3$,
is not any $f$ of Theorem~\ref{thm:fam}, and for polynomial $f$ is an
Artin--Schreier--Witt curve of a polynomial Witt vector of $s_2$-degree $\ge8$;
$D(U_3)\cap\{1,2\}=\emptyset$, $D(U_3)\supseteq\langle4,6,11\rangle\smallsetminus\{0\}$,
and $D(U_3)\cap\{3,5,7,9,13\}$ is unknown. \stag{proved}
\end{theorem}

The open questions of Section~9.2 are refined as follows. (Q1): is there
$f\in K/\wp K$, with at least two poles on the $s_2$-line for general $(s_0,s_1)$ or a
polynomial of $s_2$-degree $\ge8$, such that $K(v_f)$ is a witness? The first shape
requires the \'etale $\Z/8$-isogeny torsor of an ordinary elliptic curve
(Appendix~\ref{app:isog}); the level-$2$ version (is $(s_0,s_1)_K\otimes K(v_f)$ the
\'etale $\Z/4$-isogeny torsor of an ordinary elliptic curve at a non-constant point?) is
the natural first step, since a negative answer there kills all such level-$3$
witnesses by truncation. (Q2): does \eqref{eq:N1} fail for all $\alpha,\beta\in K$
(a local analysis at $s_2=\infty$ for pole orders $\ge2$ is the missing piece), and
which of $3,5,7,9$ lie in $D(U_3)$? The splitting type over $K$ of the degree-$11$
fibres of the pole-order-$11$ pencil is a finite computation that would decide the
latter. (Q3), (Q4) are unchanged, except that $d_1^{(p')}(\tau_2)\le p^2-p+1$ for all
$p$ and $=3$ for $p=2$.

\subsection{A quantitative Reichstein--Vistoli theorem}\label{ss:qRV}

Reichstein--Vistoli \cite{RV} prove $\ed_{k,p}(G)\le1$ for every finite $p$-group $G$ in
characteristic $p$, i.e.\ every $G$-torsor becomes one-dimensional over \emph{some}
extension of degree prime to $p$: in the language of the profile, $d_1^{(p')}(\tau)<\infty$
for every $\tau$. The proof passes through the $p$-closure, a large field, and gives
neither a bound on the degree nor uniformity in $\tau$. The Riemann--Roch argument of
Theorem~\ref{thm:B} gives both. Throughout this subsection $k$ is a field of
characteristic $p$, algebraically closed unless said otherwise, $G$ a finite group, and
$V$ a faithful representation of $G$ over $k$ with generic torsor $\tau_\gen$ over
$K=k(V)^G$.

\begin{proposition}[the generic torsor dominates every profile]\label{prop:vers}
Let $\tau$ be a $G$-torsor over an infinite field $F\supseteq k$. Then for every $d\ge1$
\[
\edd d_k(\tau)\le\edd d_k(\tau_\gen),\qquad
\edd d_k{}'(\tau)\le\edd d_k{}'(\tau_\gen)
\]
(the second for the prime-to-$p$ profiles), hence $d_j(\tau)\le d_j(\tau_\gen)$ and
$d_j^{(p')}(\tau)\le d_j^{(p')}(\tau_\gen)$ for all $j$. Consequently
$\sup_\tau\edd d_k(\tau)=\edd d_k(\tau_\gen)$, and the profile of $\tau_\gen$ does not
depend on the faithful representation $V$. \stag{proved}
\end{proposition}

\begin{proof}
Let $F'/K$ realise $\edd d_k(\tau_\gen)=m$, $[F':K]=d'\le d$; by Proposition~3.6 (or
Lemma~\ref{lem:pi}) we may take $F'/K$ separable. By Theorem~\ref{thm:geom} there are a
quasi-projective generically free primitive $G$-variety $Z$ of dimension $\le m$ and a
$G$-equivariant morphism $\operatorname{Spec}(L\otimes_KF')\to Z^{\rm free}$. As in the
proof of Theorem~4.2, let $W'$ be the normalisation of $V$ in $L\otimes_KF'$; then
$\pi\colon W'\to V$ is finite, $G$-equivariant, generically \'etale of degree $d'$, and
the equivariant map $W'\dashrightarrow Z^{\rm free}$ is defined on a dense $G$-stable
open $U$. Choose a dense $G$-stable open $V_1\subseteq V$ over which $\pi$ is \'etale and
$\pi^{-1}(V_1)\subseteq U$. Twisting by $\tau$ (functoriality of twists,
\cite[\S3]{DR}) gives a finite \'etale $F$-morphism ${}^\tau\pi^{-1}(V_1)\to{}^\tau V_1$
of degree $d'$ and an $F$-morphism ${}^\tau\pi^{-1}(V_1)\to{}^\tau Z^{\rm free}$. Now
${}^\tau V$ is an $F$-form of the vector space $V_F$, hence isomorphic to affine space by
Hilbert~90, so ${}^\tau V_1(F)$ is non-empty as $F$ is infinite; pick $x\in{}^\tau V_1(F)$.
Its fibre is $\operatorname{Spec}A$ for an \'etale $F$-algebra $A=\prod_iF_i$ of
dimension $d'$, and each $\operatorname{Spec}F_i\to{}^\tau Z^{\rm free}$ gives a closed
point of ${}^\tau Z$ in the free locus of degree $\le[F_i:F]$. The degrees $[F_i:F]$
sum to $d'\le d$, so $\edd d_k(\tau)\le m$ by Theorem~\ref{thm:geom}; if $p\nmid d'$,
some $[F_i:F]$ is prime to $p$. The last two statements: the supremum is attained at
$\tau_\gen$, and two generic torsors (for $V$ and $V'$) are torsors over infinite
fields, so each dominates the other.
\end{proof}

\begin{definition}
For a smooth projective curve $C$ over $k$ with a faithful $G$-action and a $G$-fixed
$k$-point $\infty$, let $\lambda(C,\infty)$ be the smallest element of the Weierstrass
semigroup $H(\infty)$ that is prime to $p$, and let $\lambda_k(G)$ be the minimum of
$\lambda(C,\infty)$ over all such pairs ($\lambda_k(G)=\infty$ if none exists).
\end{definition}

Such pairs exist for every finite $p$-group and, more generally, for every
$G\cong P\rtimes\Z/m$ with $P$ a $p$-group and $p\nmid m$: by Harbater and Katz--Gabber
\cite{Ha,Ka} there is a connected $G$-cover of $\mathbb P^1$ \'etale over
$\mathbb G_m$ (over $\A^1$ if $G$ is a $p$-group) with inertia group $G$ at $\infty$, and
the point above $\infty$ is then $G$-fixed. Conversely a $G$-fixed point forces $G$ to be
an inertia group in characteristic $p$, hence of this form. For $G=\Z/p^n$ the pair
$(\bar U_n,\infty)$ is defined over $\mathbb F_p$.

\begin{theorem}[quantitative Reichstein--Vistoli]\label{thm:qRV}
\begin{enumerate}
\item[(a)] Let $(C,\infty)$ be as in the definition. For every $G$-torsor $\tau$ over
every field $F\supseteq k$ there is a separable extension $F'/F$ with $p\nmid[F':F]\le
\lambda(C,\infty)$ and $\ed_k(\tau\otimes F')\le1$. Hence
\[
d_1^{(p')}(\tau)\le d_1^{(p')}(\tau_\gen)\le\lambda_k(G)\qquad\text{for every }\tau .
\]
\item[(b)] For $G=\Z/p^n$: $\lambda_k(\Z/p^n)\le l_n(p)=1+(p-1)\sum_{j=1}^{n-1}p^{2j-1}$,
realised by $(\bar U_n,\infty)$, and $\lambda(C,\infty)=l_n(p)$ is the minimum of
$\lambda(C,\infty)$ over all connected $\Z/p^n$-covers $C\to\mathbb P^1$ \'etale over
$\A^1$. So every $\Z/p^n$-torsor over every $F\supseteq k$ becomes one-dimensional over
a prime-to-$p$ extension of degree at most $l_n(p)$; the bound is $1$ for $n=1$,
$p^2-p+1$ for $n=2$, $11$ at $(2,3)$, $43$ at $(2,4)$, and it is attained by
$\tau_\gen$ for $p^n=4$ (Theorem~\ref{thm:A}). Together with Proposition~3.4,
$d_1(\tau)\le\min(p^{n-1},l_n(p))$ for every $\Z/p^n$-torsor.
\item[(c)] For an arbitrary finite $G$ with Sylow $p$-subgroup $G_p$:
$d_1^{(p')}(\tau)\le[G:G_p]\cdot\lambda_k(G_p)$ for every $G$-torsor $\tau$ over every
$F\supseteq k$.
\end{enumerate}
\stag{proved; (b) uses the inequality $u_{i+1}\ge pu_i$ for upper breaks of cyclic
$p^n$-extensions of $k((x))$ \cite{Ga,Th} and the existence results \cite{Ha,Ka}, both
cited}
\end{theorem}

\begin{proof}
(a) If $F$ is finite, $\ed_k(\tau)=0$. Let $F$ be infinite, $X={}^\tau C$, $\infty'\in X(F)$
the twist of $\infty$. The non-free locus $N\subset C$ is a finite $G$-stable set (fixed
points of the nontrivial elements of a group acting faithfully on a curve), so ${}^\tau N$
is a finite set of closed points of $X$. The Weierstrass semigroup of $\infty'$ on $X$
equals $H(\infty)$ (Riemann--Roch dimensions are invariant under base change, and
$(X,\infty')_{\bar F}\cong(C,\infty)_{\bar F}$), so there is $f\in F(X)$ with pole divisor
exactly $\lambda\infty'$, $\lambda=\lambda(C,\infty)$. The fibres $f=c$, $c\in F$, are
pairwise disjoint effective $F$-rational divisors of degree $\lambda$ not containing
$\infty'$; all but finitely many avoid ${}^\tau N$. For such a fibre the degrees of its
closed points sum to $\lambda$, which is prime to $p$, so one closed point has degree
$e$ prime to $p$, $e\le\lambda$, and lies in the free locus; its residue field $F'$ is
separable over $F$ (its degree is prime to $p$), and Theorem~\ref{thm:geom} with
$Z=C^{\rm free}$ gives $\ed_k(\tau\otimes F')\le1$. The chain of inequalities is
Proposition~\ref{prop:vers}.

(b) The upper bound is Theorem~\ref{thm:B}. For a connected $\Z/p^n$-cover
$C\to\mathbb P^1$ \'etale over $\A^1$ the argument in the proof of Theorem~\ref{thm:B}
applies verbatim ($C\to C/\langle g^{p^{n-1}}\rangle$ is an Artin--Schreier cover \'etale
over the affine part, so the coordinate ring of $C\smallsetminus\{\infty\}$ is free with
basis $1,y,\dots,y^{p-1}$), and shows $\lambda(C,\infty)=l_n(C,\infty)$, the last lower
ramification break at $\infty$. Writing the upper breaks as $u_1$ and
$u_{j+1}=pu_j+\delta_j$ with $\delta_j\ge0$ \cite{Ga,Th}, Herbrand's formula gives
$l_n=u_1+\sum_{j=1}^{n-1}p^j\bigl((p-1)u_j+\delta_j\bigr)$, which is minimal for
$u_1=1$ and all $\delta_j=0$, i.e.\ for the upper breaks $(1,p,\dots,p^{n-1})$ of
$\bar U_n$ (Appendix~\ref{app:genus}). Sharpness at $p^n=4$ is Theorem~\ref{thm:A}; the
last claim is Proposition~3.4 applied to an arbitrary torsor ($\edd{p^{n-1}}_k(\tau_a)\le
\tr_kk(a_0)\le1$).

(c) $T/G_p\to\operatorname{Spec}F$ is \'etale of degree $[G:G_p]$, so it has a field
factor $F_1$ with $p\nmid[F_1:F]\le[G:G_p]$; over $F_1$ the torsor has a point of
$T/G_p$, hence $\tau\otimes F_1\cong\Ind_{G_p}^G\tau_p$ for a $G_p$-torsor $\tau_p$ over
$F_1$. By (a) there is $F_2/F_1$ with $p\nmid[F_2:F_1]\le\lambda_k(G_p)$ and
$\ed_k(\tau_p\otimes F_2)\le1$, and $\ed_k(\tau\otimes F_2)=\ed_k(\tau_p\otimes F_2)$ by
Lemma~\ref{lem:ind}.
\end{proof}

\begin{remark}
(i) What is new relative to \cite{RV} and to Proposition~3.7(b),(c): the explicit bound,
its uniformity in $\tau$, the identification of the optimal bound from one-point covers
with a local ramification invariant for cyclic groups, and the versality of the profile
(Proposition~\ref{prop:vers}), which also answers the question whether the profile of
$\tau_\gen$ depends on $V$. The mechanism, a free point on a curve with a fixed point, is
that of \cite[Lemma~9.1]{RS} and Proposition~3.7(c); what replaces largeness is
Riemann--Roch on the twisted curve. (ii) A crude bound valid for every $(C,\infty)$ is
$\lambda(C,\infty)\le$ the least integer prime to $p$ that is $\ge2g(C)$, since
$H(\infty)\supseteq[2g(C),\infty)$; at $(2,3)$ this gives $15$ against $11$.
(iii) Whether $\lambda_k(\Z/p^n)=l_n(p)$ over \emph{all} curves with a fixed point, not
only one-point covers, is open, as is the exact value of $d_1^{(p')}(\tau_\gen)$ for
$n=2$ and $p$ odd (between $2$ and $p^2-p+1$) and at $(2,3)$ (in $\{3,5,7,9,11\}$).
(iv) The uniform statement is a statement about $\sup_\tau$: by
Proposition~\ref{prop:vers} it is equivalent to the bound for $\tau_\gen$ alone.
\end{remark}
